\section{点态极限为何不能直接穿过积分号}
\label{sec:12-Real-Analysis-Support/04-Limits-and-Integration/01}

你手头有一列函数 \(f_n\)，各自可积。它们的逐点极限很漂亮——每个 \(x\) 处 \(f_n(x)\) 的趋向都清清楚楚。直觉告诉你：
\[
\lim_{n\to\infty}\int f_n \stackrel{?}{=} \int \lim_{n\to\infty}f_n.
\]
把极限号穿过积分号，先做逐点极限再积分，往往让被积函数大为简化。这一步若总是合法，许多估计可以少写好几页。

在 \([0,1]\) 上试试下面这列函数：
\[
f_n(x)=n\,\mathbf1_{(0,1/n)}(x).
\]
对每个固定 \(x>0\)，当 \(n\) 足够大时 \(1/n<x\)，指示函数归零，所以 \(f_n(x)=0\)；\(x=0\) 处也定义为零。于是逐点处处有
\[
f_n(x)\to 0.
\]
再算每个 \(f_n\) 的积分：
\[
\int_0^1 f_n(x)\,dx = n\cdot\frac1n = 1.
\]
左边：\(\lim_n \int f_n = 1\)。右边：\(\int \lim_n f_n = \int 0 = 0\)。不等。

发生了什么？尖峰越来越高（高度 \(n\)），越来越窄（宽度 \(1/n\)），向零点滑去。逐个固定点看，尖峰最终经过并消失——``经过''这一事实不体现在点态极限里。总面积却始终维持在一。点态收敛只记录每个\emph{位置}的最终去向，不记录质量在空间中的移动和集中。

\subsection{统一误差控制：一致收敛直接解决}

如果想在 Riemann 积分框架下越过这道障碍，最直接的条件是一致收敛。若 \(f_n\to f\) 在 \([a,b]\) 上一致，且各 \(f_n\) 可积，则
\[
\abs{\int_a^b f_n - \int_a^b f}
\le \int_a^b |f_n-f|
\le (b-a)\,\norm{f_n-f}_\infty,
\]
其中 \(\norm{\cdot}_\infty\) 是上确界范数，一致收敛正是说它趋于零。右端归零给出积分的交换。

注意``有限区间''是前提的一部分。若积分域总长度无限，统一误差上界 \(\varepsilon_n\to0\) 乘上无限体积，积分差仍无法控制。

一致收敛是足够强的条件——强到排除了所有质量逃逸的可能——却不是必需的。\(f_n(x)=x^n\) 在 \([0,1]\) 上的逐点极限在 \(x=1\) 处间断，所以收敛不一致，但积分 \(\int_0^1 x^n dx = 1/(n+1)\to 0\)，极限仍穿过积分号。这说明还存在更精细的工具。

\subsection{只向上增长时，质量不会凭空丢失}

如果函数非负，且在整个过程中只增加不减少，交换就自动成立。

\begin{theorem}[单调收敛定理，MCT]
若 \((f_n)\) 是非负可测函数，且
\[
f_n(x)\uparrow f(x)
\]
几乎处处（即每点处单调递增到 \(f\)），则
\[
\int f_n \uparrow \int f,
\]
允许两侧为 \(+\infty\)。
\end{theorem}

例如在 \((0,1]\) 上定义
\[
f_n = \mathbf1_{[1/n,1]}.
\]
集合随 \(n\) 增大而向右扩张，\(f_n\uparrow 1\)，积分 \(1-1/n \uparrow 1\)。

非负性堵住了正负抵消的通道——不能这边涨那边跌把总面积稳住。单调性则保证旧质量一旦进入就不会在下一步消失或迁走。它不要求统一可积上界，也不要求一致收敛，两个条件精确对应了尖峰逃逸的两种机制：符号翻转和质量漂移。

若函数可能为负，可以通过加上一个可积下界转化为非负情形，或使用下一节的控制收敛定理。

\subsection{即使不等号，也要给一个下界：Fatou 引理}

很多时候序列并非单调——这时交换不能保证，但尚可控制质量丢失的\emph{方向}。

\begin{theorem}[Fatou 引理]
对非负可测函数 \((f_n)\)，
\[
\int \liminf_{n\to\infty} f_n \;\le\; \liminf_{n\to\infty} \int f_n.
\]
\end{theorem}

翻译成画面：极限函数的积分绝不会超过积分序列的下极限。尖峰例子中左侧 0、右侧 1，严格不等——质量在极限过程中弄丢了，Fatou 用不等号记录了丢失方向。

Fatou 的典型应用场景是变分法和概率论中的下半连续性论证：近似解的能量在极限中不会突然比 liminf 更大。统计风险分析中，``渐近风险 \(\ge\) 极限风险''用的也是同一不等式骨架。

单调收敛定理可以从 Fatou 推出：若 \(f_n\uparrow f\)，Fatou 给出 \(\int f\le\liminf\int f_n\)；而 \(f_n\le f\) 给出 \(\limsup\int f_n\le\int f\)；两侧夹拢即等号。反过来，Fatou 的证明思路是把序列换成尾部下确界 \(g_n=\inf_{k\ge n}f_k\)——它单调上升到 \(\liminf f_n\)——再对 \(g_n\) 用 MCT。两个定理互为表里。

\subsection{动手前先排查质量可能从哪条路逃逸}

在面对
\[
\lim_n\int f_n \stackrel{?}{=} \int \lim_n f_n
\]
时，问自己这几个问题比直接计算更有效：

\begin{itemize}
\tightlist
\item 是否只有点态收敛？——点态收敛不约束质量移动，仅此一点就不足以保证交换。
\item 高度是否可能变大形成尖峰？——如果是，MCT 或 DCT（下一节）是出路。
\item 质量是否可能向无穷远漂移？——积分域非紧时尤其常见。
\item 正负部分是否分别可积，有没有因符号抵消而掩盖总质量变化？
\item 是否单调非负？——MCT 直接接管。
\item 是否存在共同可积包络？——下一节的 DCT 接管。
\item 积分域是否随 \(n\) 变化？——需要 Leibniz 积分法则处理。
\end{itemize}

这些问题决定了结论，而``函数看起来都趋于零''不说明任何事——移动尖峰恰恰是看起来趋于零、积分却稳如磐石的最强反例。交换之前先定位可能的逃逸通道，是整条实分析积分极限路径的基本操作习惯。
